\chapter{混音总览}
\label{chap:mixing-overview}

混音不是把每一轨都调大，而是让听众在整首歌里听清该听的东西：主唱、节奏、
低频与空间各在其位。本章先给夜半的信号路径、监听入口和一套可执行的混音顺序；
具体参数与效果器在后续章节逐项展开。

\section{混音的目标}

一首歌的混音至少要同时满足三件事：

\begin{itemize}
  \item \term{可懂}：主唱与主要旋律清楚，不被鼓、贝斯或混响盖住。
  \item \term{平衡}：低频、中频、高频没有明显缺口或堆积，整体不靠某一轨硬撑。
  \item \term{可控}：从大到小、从静到响都有余量，导出后在不同设备上不破音、不塌陷。
\end{itemize}

“好听”是审美判断，“可懂、平衡、可控”是可以逐项检查的工程目标。先满足后者，
再谈风格。

\section{夜半的信号路径}

在夜半里，每条轨道先经过自己的设备链，再经过推子与声像，汇入辅助总线或主总线；
主总线再把最终信号送到监听控制室。设备链的顺序就是信号流顺序：乐器在前，
效果器在后。

\begin{center}
\begin{tikzpicture}[
  node distance=8mm and 10mm,
  box/.style={draw, rounded corners=2pt, minimum width=27mm, minimum height=8mm, align=center, font=\small},
  bus/.style={draw, rounded corners=2pt, minimum width=27mm, minimum height=8mm, align=center, font=\small, fill=black!4},
  arrow/.style={-{Latex[length=2mm]}, thick}
]
  \node[box] (tracks) {各轨设备链\\乐器 / 效果};
  \node[box, right=of tracks] (strip) {通道条\\推子 + 声像};
  \node[bus, right=of strip] (aux) {辅助总线\\共用效果};
  \node[bus, right=of aux] (master) {主总线\\Master};
  \node[box, right=of master] (room) {监听控制室\\电平 / 响度};

  \draw[arrow] (tracks) -- (strip);
  \draw[arrow] (strip) -- (aux);
  \draw[arrow] (aux) -- (master);
  \draw[arrow] (master) -- (room);
\end{tikzpicture}
\end{center}

混音控制台（X / F8）把每条轨显示成一条通道条：推子、声像、静音/独奏，以及
这条轨的设备链入口。左侧检查器（I）显示同一条轨的通道条细节，与混音台是同一份
状态；在一边改动，另一边同步。所有轨道最终汇入主总线。

\section{监听：先相信读数，再相信耳朵}

监听控制室（F4，或“混音 → 显隐 监听控制室”）是独立的主控监听区。它显示主总线的
积分响度、短时响度与峰值读数，并提供 DIM（−20 dB）、MONO、SIDE（S=L−R）等试听开关。这些开关只改变
监听方式，不动任何一条轨道的推子；主总线推子仍是工程状态，改动会被保存。

\screenshot{mixer-console}{混音控制台打开，至少三条轨有通道条，主总线推子与电平可见。}

\screenshot{control-room}{监听控制室展开，积分/短时 LUFS 与峰值 dBTP 读数可见。}

监听环境本身也会骗人。用耳机混音容易把低频做得过重，用客厅音箱混音容易把
高频做得过亮。至少做两次检查：一次在常用监听上，一次换一副耳机或换一个音量；
两次都成立的平衡，才更可能是真的平衡。

\section{参考曲目}

参考曲目不是拿来照抄旋律，而是用来校准“这首歌应该有多响、多亮、多宽”。做法：

\begin{enumerate}
  \item 用“文件 → 导入音频文件...”把参考曲导入素材池。
  \item 在一条单独的音频轨上摆放参考曲，不要把它和成品轨叠在同一轨。
  \item 旁通这条参考轨上的效果器，让参考曲尽量接近原始素材的状态播放。
  \item 在控制室看参考曲的积分响度与峰值，记住这个量级；切回自己的工程，比较
        同样的段落（例如副歌对副歌），不要拿整首对一段。
\end{enumerate}

参考曲只用来做 A/B 判断，不进入最终导出。导出前把参考轨静音或删除，避免它被
一并渲染。

\section{一套可执行的混音顺序}

下面的顺序不是硬性规定，但它能避免“先加混响、最后发现电平全错”的返工：

\begin{enumerate}
  \item \term{粗平衡}：所有轨推子先归到合理位置，不削不爆；用独奏逐轨确认没有
        误静音、误独奏或漏摆的片段。
  \item \term{增益结构}：先让每条轨在自身设备链里有合适的工作电平，再用推子做
        相对平衡。推子负责“谁更响”，不负责“修失真”。
  \item \term{EQ}：先修掉掩蔽与共振，再谈音色。通常先处理低频冲突与中频掩蔽，
        高频留到最后微调。
  \item \term{动态}：压缩、限制、门限按需要上；先问“要解决什么问题”，再选效果器，
        不要因为插件栏里有它就加。
  \item \term{空间}：混响与延迟优先放在辅助总线上多轨共用，统一空间感；近距离
        效果再考虑插入式。
  \item \term{自动化}：段落之间的音量、声像、发送与效果参数用自动化写清楚；
        不要在段落交界处突然手动改值。
  \item \term{总线与收尾}：看控制室的响度与峰值，做总线处理与导出前检查；
        导出后再听一遍成品，不只看工程里的表头。
\end{enumerate}

每一步只解决一类问题。若一步之后更差，先撤回这一步，再决定是否换方法；
不要在上一步没解决时叠加更多效果器。

\section{练习}

\begin{enumerate}
  \item 选一段 8 小节，按“粗平衡 → 增益结构 → EQ → 动态 → 空间 → 自动化 → 总线收尾”走一遍；每一步只记一个决定。
  \item 用同一段参考曲在控制室做正常、DIM、MONO、SIDE 四种监听；写下至少三条差别，不要用参考曲的响度直接判断自己的电平。
  \item 把这次走查的七步顺序写成工程备注；如果某一步没有做，写明原因，不要留空。
\end{enumerate}

\section{本章与帮助主题的映射}

\begin{table}[htbp]
\centering
\small
\begin{tabular}{@{}p{0.34\linewidth}p{0.52\linewidth}@{}}
\hline
本章小节 & 应用内帮助锚点 \\
\hline
混音的目标 & \texttt{mixing-goals} \\
夜半的信号路径 & \texttt{mixing-signal-flow} \\
监听：先相信读数，再相信耳朵 & \texttt{mixing-monitoring} \\
参考曲目 & \texttt{mixing-reference-track} \\
一套可执行的混音顺序 & \texttt{mixing-workflow} \\
\hline
\end{tabular}
\caption{第一章帮助主题映射}
\end{table}
